\documentclass[../main.tex]{subfiles}

\begin{document}
    \subsection{Bagi THR}

\begin{code}
\begin{verbatim}
#include <iostream>
\end{verbatim}
\end{code}

    {\Script} di atas digunakan untuk membolehkan penggunaan fungsi-fungsi dan objek-objek yang berada di dalam {\header} {\file} \citecode{iostream}. {\Header} {\file} \citecode{iostream} berisi hal-hal yang dibutuhkan untuk membantu proses {\ioi} dan {\ioo}.

\begin{code}
\begin{verbatim}
using namespace std;
\end{verbatim}
\end{code}

    {\Script} \citecode{using namespace std;} adalah {\statement} yang digunakan untuk memberikan perintah kepada {\compiler} bahwa file tersebut akan menggunakan seluruh fungsi yang menjadi bagian dari \citecode{namespace std}.

\begin{code}
\begin{verbatim}
int main() 
\end{verbatim}
\end{code}

    {\Script} digunakan sebagai \textit{entry point} program tersebut yang akan dijalankan setiap kali program Bagi THR dijalankan. \textit{Entry point} diperlukan oleh setiap program yang ditulis pada bahasa C++ untuk berjalan.

\begin{code}
\begin{verbatim}
string name = " ";
cout << "Masukan Nama  : ";
getline(cin, name);
\end{verbatim}
\end{code}

    {\Script} di atas digunakan untuk mendapatkan masukan dari pengguna yang nantinya akan dimasukan ke dalam variabel \citecode{name} menggunakan fungsi \citecode{getline}. Fungsi \citecode{getline} digunakan untuk mendapatkan masukan yang bisa berisikan karakter spasi.

\begin{code}
\begin{verbatim}
int type = 0;
cout << "Yoo!!! " << name << ", Kamu Tipe Mahasiswa Seperti Apa?\n";
cout << "1. Mahasiswa Kupu-kupu\n";
cout << "2. Mahasiswa Kura-kura\n";
cout << "3. Mahasiswa Kuda-kuda\n";
while (true) {
  cout << "Masukan Pilihan Kamu: ";
  cin >> type;
  if (cin.fail() || (type < 1 || type > 3)) {
    cin.clear();
    cin.ignore();
    cout << "Input tidak valid!" << endl;
    continue;
  } break;
}
\end{verbatim}
\end{code}

    {\Script} di atas digunakan untuk menunjukan pilihan dan mendapatkan masukan yang dapat bertipe pilihan tersebut dari pengguna yang akan disimpan pada variabel \citecode{type}. Masukan akan terus diminta apabila masukan tidak sesuai dengan masukan yang diminta yakni masukan berupa angka yang lebih dari 1 dan kurang dari 3.

\begin{code}
\begin{verbatim}
int money = 0;
switch (type) {
case 1:
  cout<<"Selamat Datang "<<name<<" - Mahasiswa Kuliah-Pulang\n";
  cout<<"Santai banget ya kamu :)\n"; money = 500000; break;
case 2:
  cout<<"Selamat Datang "<<name<<" - Mahasiswa Kuliah-Rapat\n";
  cout<<"Kasian banget ya kamu :)\n"; money = 850000; break;
case 3:
  cout<<"Selamat Datang "<<name<<" - Mahasiswa Kuliah-Dagang\n";
  cout<<"Semangat ya buat kamu :)\n"; money = 1000000; break; }
\end{verbatim}
\end{code}

    {\Script} di atas digunakan untuk menentukan nilai variabel \citecode{money} sesuai nilai variabel \citecode{type}. Keluaran berupa teks juga akan ditunjukan akan berbeda sesuai dengan nilai variabel \citecode{type}. {\Statement} \citecode{switch} digunakan karena terbatasnya variasi nilai variabel \citecode{type}.

\begin{code}
\begin{verbatim}
int duration = 0;
while (true) {
  cout << "Berapa Jam waktu Kuliahmu Ramadhan ini? ";
  cin >> duration;
  if (cin.fail()) {
      cin.clear();
      cin.ignore();
      cout << "Input tidak valid!" << endl;
      continue;
  } break;
}
\end{verbatim}
\end{code}

    {\Script} di atas digunakan untuk menentukan nilai variabel \citecode{duration} yang akan menyimpan durasi mahasiswa kuliah. Variabel \citecode{duration} bertipe data \citecode{int} yang digunakan untuk menyimpan bilangan bulat.

\begin{code}
\begin{verbatim}
cout<<"------------------------------------------\n";
cout<<"THR buat kamu sebesar   : Rp."<< money << "\n";
cout<<"Durasi Kuliah Minggu ini: "<<(duration *= 3600)<<" Detik\n";
cout<<"------------------------------------------\n";
\end{verbatim}
\end{code}

    {\Script} di atas digunakan untuk menunjukan nilai variabel \citecode{money} dan nilai variabel \citecode{duration} kepada pengguna. Kedua nilai tersebut bertipe data \citecode{int} yang digunakan untuk menyimpan bilangan bulat.

\begin{code}
\begin{verbatim}
char overtime = 'y';
while (true) {
  cout << "Apakah Kamu Pernah Kuliah Sampai Malam? (Y/N): ";
  cin >> overtime;
  if (cin.fail() ||
      (overtime != 'Y' && overtime != 'N')
      &&(overtime!='y'&&overtime!='n')){cin.clear();cin.ignore();
      cout << "Input tidak valid!" << endl; continue; } break; }
if (overtime == 'Y' || overtime == 'y') { int duration1 = 0;
  cout << "Kamu Memang Rajin !!!\n";
  cout << "Holil Ingin Memberimu Lebih Banyak THR\n";
  while (true) {cout<<"Masukan Total Waktu Kuliah Malam Kamu (Detik): ";
    cin >> duration1;
    if (cin.fail()) { cin.clear(); cin.ignore();
      cout << "Input tidak valid!" << endl; continue; } break; }
  duration += duration1;
  if ((duration1/3600.0f)>1.0f&&(duration1/3600.0f)<=2.0f){
    money += 200000;
  }else if((duration1/3600.0f)>2.0f&&(duration1/3600.0f)<=3.0f){
    money += 500000;
  }else if((duration1/3600.0f)>3.0f) {
    money *= 2;
  }
}
\end{verbatim}
\end{code}

{\Script} di atas digunakan untuk menanyakan apakah mahasiswa kuliah malam. Bila iya akan ditanyakan berapa total waktu kuliah malamnya kemudian akan ditambahkan uang ke dalam variabel \citecode{money} sesuai waktu kuliah malamnya.

\begin{code}
\begin{verbatim}
cout<<"Durasi Kuliah Ramadhan Ini: "<<(duration/3600)<<" Jam,"
    <<(duration/60)-((duration/3600)*60)<<" Menit,"
    <<(duration-((duration/3600)*3600))-
    (((duration-((duration/3600)*3600))/60)*60)<<" Detik\n";
\end{verbatim}
\end{code}

{\Script} di atas digunakan untuk menampilkan durasi kuliah dari nilai variabel \citecode{duration} sesuai format jam, menit dan detik. Untuk mendapat menit dan detik dilakukan {\type} \textit{punning} dengan membagi \citecode{duration} dengan 3600 dalam kurung dan mengalikannya kembali dengan 3600 untuk menghilangkan angka berlebihan.

\subsection{Motif Baju}

\begin{code}
\begin{verbatim}
#include <conio.h> 
\end{verbatim}
\end{code}

{\Script} di atas digunakan untuk membolehkan penggunaan fungsi-fungsi dan objek-objek yang berada di dalam {\header} {\file} \citecode{conio.h}. {\Header} {\file} \citecode{conio.h} berisi hal-hal yang dibutuhkan untuk berinteraksi dengan konsol sistem operasi Microsoft Windows dan DOS.

\begin{code}
\begin{verbatim}
#include <iostream>
\end{verbatim}
\end{code}

{\Script} di atas digunakan untuk membolehkan penggunaan fungsi-fungsi dan objek-objek yang berada di dalam {\header} {\file} \citecode{iostream}. {\Header} {\file} \citecode{iostream} berisi hal-hal yang dibutuhkan untuk membantu proses {\ioi} dan {\ioo}.

\begin{code}
\begin{verbatim}
using namespace std;
\end{verbatim}
\end{code}

{\Script} \citecode{using namespace std;} adalah {\statement} yang digunakan untuk memberikan perintah kepada {\compiler} bahwa file tersebut akan menggunakan seluruh fungsi yang menjadi bagian dari \citecode{namespace std}.

\begin{code}
\begin{verbatim}
int main() 
\end{verbatim}
\end{code}

{\Script} digunakan sebagai \textit{entry point} program tersebut yang akan dijalankan setiap kali program Bagi THR dijalankan. \textit{Entry point} diperlukan oleh setiap program yang ditulis pada bahasa C++ untuk bisa berjalan dengan normal.

\begin{code}
\begin{verbatim}
char state = 0; int n0 = 0; int n1 = 0;
\end{verbatim}
\end{code}

{\Script} di atas digunakan untuk menciptakan tiga buah variabel yakni \citecode{state} yang akan digunakan untuk menyimpan keadaan terkini program dan \citecode{n0} serta \citecode{n1} yang digunakan untuk menyimpan masukan atau perhitungan sementara.

\begin{code}
\begin{verbatim}
while (true) {
  cout << "\e[H\e[2J\e[0m";
\end{verbatim}
\end{code}

{\Script} di atas merupakan awal perulangan program Motif Baju. Perulangan tersebut akan berulang selamanya setelah program dijalankan. Baris \citecode{cout << "\textbackslash e[H\textbackslash [2J\textbackslash [0m";} digunakan untuk membersihkan isi layar program dengan menggunakan karakter ANSI.

\begin{code}
\begin{verbatim}
switch (state & 0x1) { case 0: {
  cout << "-----------------------------------" << endl;
  cout << "|          MOTIF BATIK            |" << endl;
  cout << "|             HOLIL               |" << endl;
  cout << "-----------------------------------" << endl;
  if ((state >> 1) == 0) cout << "\e[35m";
  cout << "-----------------------------------" << endl;
  cout << "|          BANGUN DATAR           |" << endl;
  cout << "-----------------------------------\e[39m" << endl;
  if ((state >> 1) == 1) cout << "\e[35m";
  cout << "-----------------------------------" << endl;
  cout << "|          POHON CEMARA           |" << endl;
  cout << "-----------------------------------\e[39m" << endl;
  if ((state >> 1) == 2) cout << "\e[35m";
  cout << "-----------------------------------" << endl;
  cout << "|          POLA BILANGAN          |" << endl;
  cout << "-----------------------------------\e[39m" << endl;
  if ((state >> 1) == 3) cout << "\e[35m";
  cout << "-----------------------------------" << endl;
  cout << "|             KELUAR              |" << endl;
  cout << "-----------------------------------\e[39m" << endl;
\end{verbatim}
\end{code}

{\Script} di atas merupakan kumpulan sebuah objek \citecode{cout} yang akan digunakan untuk menampilkan sebuah menu pada pengguna. Opsi yang dipilih akan diwarnai oleh warna ungu menggunakan kode karakter ANSI. Karakter ANSI \citecode{\textbackslash e[39m} memberi warna ungu.

\begin{code}
\begin{verbatim}
n0 = getch();
if (n0 == 13) {
  state |= 0x1;
}
\end{verbatim}
\end{code}

{\Script} di atas digunakan untuk mendeteksi apakah tombol {\enter} sudah ditekan. Pendeteksian tersebut dilakukan dengan menggunakan fungsi \citecode{getch()} yang tersedia pada {\header} {\file} \citecode{conio.h}.

\begin{code}
\begin{verbatim}
if (n0 == 72)
  state = (((state >> 1) - 1) << 1);
else if (n0 == 80)
  state = (((state >> 1) + 1) << 1);
if ((state >> 1) > 3)
  state = 0x0;
else if ((state >> 1) < 0) state = 0x6;
\end{verbatim}
\end{code}

{\Script} di atas digunakan untuk mendeteksi apakah tombol panah atas atau tombol panah bawah sudah ditekan. Bila tombol panah atas atau tombol panah bawah ditekan maka nilai bit selain bit pertama \citecode{state} akan naik atau turun sesuai kondisi.

\begin{code}
\begin{verbatim}
case 1: { n0 = 0; n1 = 0;
  switch ((state >> 1)) {
\end{verbatim}
\end{code}

{\Script} di atas digunakan untuk mengganti tampilan program sesuai nilai \citecode{state >> 1} atau nilai \citecode{state} selain bit pertama. Kedua \citecode{n0} dan \citecode{n1} di atur ke 0 untuk membersihkan isi dari proses sebelumnya.

\begin{code}
\begin{verbatim}
case 0: { // Bangun Datar
  cout << "-----------------------------------" << endl;
  cout << "|          BANGUN DATAR           |" << endl;
  cout << "-----------------------------------" << endl;
\end{verbatim}
\end{code}

{\Script} di atas digunakan untuk menampilkan judul pilihan yakni pilihan BANGUN DATAR. Dengan menggunakan objek \citecode{cout} dari {\header} {\file} \citecode{iostream} yang merupakan bagian dari C++ \textit{standard library} judul pilihan dapat ditampilkan.

\begin{code}
\begin{verbatim}
cout << "           Pola Persegi :)" << endl;
cout << "Masukan Jumlah Baris: "; cin >> n0;
for (int i = 0; i < n0; ++i) {
	char ch = 'A';
	for (int j = 0; j < n0; ++j) {
		if (i > 0 && i == j) { n1++; ch = (char)(65 + n1);
		} cout << ch << " ";
		if ((n1 - ((int)(ch) - 65)) > 0) ch++;
	} cout << endl; }
\end{verbatim}
\end{code}

{\Script} di atas digunakan untuk menggambar persegi yang disusun oleh alfabet yang akan urutan alphabetnya menggingkat semakin jauh titik tersebut dengan pojok kiri atas persegi. Variabel \citecode{n0} merupakan masukan panjang dan lebar persegi tersebut.

\begin{code}
\begin{verbatim}
cout << "           Pola Segitiga :|" << endl;
cout << "Masukan Jumlah Baris: "; cin >> n0;
n1 = !((n0 + 2) % 2 == 0);
for (int i = 0; i < n0; ++i) {
	for (int j  = 0; j < (n0 - i) - 1; ++j) {
	  cout << " "; }
	for (int j = 0; j < (i * 2) + 1; j += 2) {
	  if ((i + 2) % 2 == 0) cout << (char)(65 + i + n1) << " ";
	  else cout << i + 1 - n1 << " ";
	} cout << endl; }
\end{verbatim}
\end{code}

{\Script} di atas digunakan untuk menggambar segitiga berpola huruf diikuti angka yang diulangi sampai nilai variabel \citecode{n0}. Variabel \citecode{n1} digunakan untuk mendeteksi apakah \citecode{n0} merupakan bilangan genap, bila \citecode{n0} merupakan genap maka 1 akan ditambahkan pada kode karakter dan satu akan dikurangi pada angka.

\begin{code}
\begin{verbatim}
cout << "           Pola Persegi Lagi :(" << endl;
cout << "Masukan Jumlah Baris: ";
cin >> n0;
for (int i = 0; i < n0; ++i) {
	for (int j = 0; j < n0; ++j) {
		int k = (i * n0) + j + 1;
		if (k/10.0f < 1.0f) cout << "0";
		cout << k << " ";
	} cout << endl; }
\end{verbatim}
\end{code}

{\Script} di atas digunakan untuk menggambar persegi berpola bilangan yang akan meningkat sesuai jaraknya dengan pojok kiri atas persegi. Angka nol akan ditambahkan pada awal bilangan bila bilangan tersebut dibawah sepuluh.

\begin{code}
\begin{verbatim}
case 1: {
  cout << "-----------------------------------" << endl;
  cout << "|          POHON CEMARA           |" << endl;
  cout << "-----------------------------------" << endl;
\end{verbatim}
\end{code}

{\Script} di atas digunakan untuk menampilkan judul pilihan yakni pilihan POHON CEMARA. Dengan menggunakan objek \citecode{cout} dari {\header} {\file} \citecode{iostream} yang merupakan bagian dari C++ \textit{standard library} judul pilihan dapat ditampilkan.

\begin{code}
\begin{verbatim}
cout << "Masukan Tinggi: ";
cin >> n0;
cout << "Masukan Lebar: ";
cin >> n1;
for (int i = 0; i < n0; ++i) {
	for (int j = 0; j < n1; ++j) {
		for (int k = 0; k < (n1 - j + 1)
			 + (n0 - i) * 2; ++k) {
			cout << " "	; }
		for (int k = 0; k < ((j * 2) + 1) + (i * 4); k += 2) {
			cout << "* "; }
		cout << endl; }
}
for (int i = 0; i < 4; ++i) {
	for (int j = 0; j < (n1 + n0) + (n0 - 2); ++j) {
		cout << " "; }
	cout << "* * * *" << endl;
}
\end{verbatim}
\end{code}

{\Script} di atas digunakan untuk menggambar pohon cemara. Terdapat dua perulangan utama yang menyusunnya, perulangan pertama digunakan untuk menggambarkan daun pohon cemara dan perulangan kedua digunakan untuk menggambarkan batan pohon cemara.

\begin{code}
\begin{verbatim}
case 2: {
  cout << "-----------------------------------" << endl;
  cout << "|          POLA BILANGAN          |" << endl;
  cout << "-----------------------------------" << endl;
\end{verbatim}
\end{code}

{\Script} di atas digunakan untuk menampilkan judul pilihan yakni pilihan POLA BILANGAN. Dengan menggunakan objek \citecode{cout} dari {\header} {\file} \citecode{iostream} yang merupakan bagian dari C++ \textit{standard library} judul pilihan dapat ditampilkan kepada pengguna.

\begin{code}
\begin{verbatim}
cout << "Masukan nilai n: ";
cin >> n0;
if ((n0 & 0x1) == 0) { cout << "<1> ";
  for (unsigned int i = 2; i < n0; ++i) {
  	bool prime = true;
  	for (unsigned int j = 2; j < i; ++j) {
  	  if (i % j == 0) prime = false;
    }
  	if (prime) { if (n0 <= 50 && (i % 10) == 3)  {
  	  	cout << "(" << i << "*n = " << (i * n0) << ") "; continue;
    }
  	  cout << "<" << i << "> "; } } cout << endl; break;
}
\end{verbatim}
\end{code}

{\Script} di atas digunakan untuk menggambar pola bilangan prima sampai dengan nilai variabel \citecode{n0}. Bila nilai variabel \citecode{n0} lebih kecil atau sama dengan 50 maka setiap bilangan yang mengandung angka 3 akan dikalikan dengan nilai variabel \citecode{n0}. {\Script} di atas hanya akan dijalankan apabila \citecode{n0} merupakan bilangan genap.

\begin{code}
\begin{verbatim}
unsigned long long n = 1;
for (unsigned int i = 0; i < n0; ++i) {
	if (i == 0) {
		cout << n; continue; }
	n *= 11;
	cout << ", " << n; }
cout << endl;
\end{verbatim}
\end{code}

{\Script} di atas digunakan untuk menggambar pola bilangan 11 pangkat 0 sampai dengan 11 pangkat nilai variabel \citecode{n0}. {\Script} di atas hanya akan dijalankan bila {\script} sebelumnya tidak dijalankan berarti nilai variabel \citecode{n0} merupakan bilangan ganjil.

\begin{code}
\begin{verbatim}
case 3: { return 0; } break; }
\end{verbatim}
\end{code}

{\Script} di atas merupakan {\statement} \textit{early return}. {\Statement} \textit{early return} berarti bahwa sebuah fungsi atau program akan dihentikan langsung ketika menemui {\statement} ini. {\Statement} ini ditandai dengan adanya {\statement} \citecode{return} sebelum dicapainya akhir blok kode.

\begin{code}
\begin{verbatim}
cout << "Tekan tombol apapun untuk kembali!" << endl;
if (getch()) { state &= 0x0; }
\end{verbatim}
\end{code}

{\Script} di atas akan mengembalikan keadaan program ke menu apabila pengguna menekan tombol apapun. Fungsi \citecode{getch()} digunakan untuk menunggu masukan pengguna, {\statement} \citecode{if} hanya mendeteksi apakah kondisi bernilai 0 untuk salah atau bukan 0 untuk benar dan setiap kode karakter pada \textit{keyboard} bernilai lebih besar dari 0.

\subsection{Hilal}
\begin{code}
\begin{verbatim}
#include <iostream>
\end{verbatim}
\end{code}

{\Script} di atas digunakan untuk membolehkan penggunaan fungsi-fungsi dan objek-objek yang berada di dalam {\header} {\file} \citecode{iostream}. {\Header} {\file} \citecode{iostream} berisi hal-hal yang dibutuhkan untuk membantu proses {\ioi} dan {\ioo}.

\begin{code}
\begin{verbatim}
using namespace std;
\end{verbatim}
\end{code}

{\Script} di atas digunakan untuk memberikan kemampuan untuk menggunakan fungsi-fungsi dan objek-objek yang tersedia dalah {\header} {\file} \citecode{iostream}. {\Script} baris \citecode{using namespace std;} memberikan kemampuan untuk menggunakan fungsi-fungsi dan objek-objek dari \textit{standard library} tanpa menulis \citecode{std::}.

\begin{code}
\begin{verbatim}
cout << "---------------------" << endl;
cout << "Program Animasi Bulan" << endl;
cout << "---------------------" << endl;
\end{verbatim}
\end{code}

{\Script} di atas digunakan untuk menampilkan {\header} atau judul program kepada pengguna. Judul program tersebut ditampilkan dengan menggunakan objek \citecode{cout} yang tersedia dalam {\header} {\file} \citecode{iostream} yang merupakan bagian dari C++ \textit{standard library}.

\begin{code}
\begin{verbatim}
int radius;
\end{verbatim}
\end{code}

{\Script} di atas merupakan {\statement} inisialisasi sebuah variabel. Dimana dalam {\script} di atas diinisialisasikan sebuah variabel \citecode{n} bertipe \citecode{int}. Kemudian variabel \citecode{n} tersebut dapat digunakan pada kode-kode setelah baris kode {\script} tersebut.

\begin{code}
\begin{verbatim}
cout << "Masukkan Radius Bulan: ";
cin >> radius;
\end{verbatim}
\end{code}

{\Script} di atas digunakan untuk meminta masukan dari pengguna. Masukan tersebut akan disimpan ke variabel \citecode{radius} yang telah diinisialisasikan sebelumnya. Variabel \citecode{radius} tersebut digunakan untuk menyimpan jari-jari bulan yang ingin digambar.

\begin{code}
\begin{verbatim}
while (true) {
  for (int arah = -(2 * radius); arah < 2 * radius; arah++) {
    for (int y = -radius; y < radius; y++) {
      for (int x = -(radius); x < (radius); x++) {
        int jarak = x * x + y * y;
        int rsquare = (radius * radius) - radius;
        int jarak2 = (x + arah) * (x + arah) + y * y;
        int rsquare2 = (radius - 1) * (radius - 1);
        int bulan1 = jarak <= rsquare;
        int bulan2 = jarak2 >= rsquare2;
        if (bulan1 && bulan2) {
          cout << setw(2) << "*";
        } else {
          cout << setw(2) << " ";}}
      cout << endl;
    }
    cout << "\33[H\33[2J";
  }
}
\end{verbatim}
\end{code}

{\Script} di atas digunakan untuk menggambar dua buah lingkaran yang akan digunakan untuk membuat animasi bulan. Ada dua buah lingkaran yang dibuat salah satu dari dua lingkaran tersebut merupakan lingkaran bulan dan lingkaran lainnya adalah lingkaran yang akan digunakan untuk menutupi lingkaran sebelumnya untuk mensimulasikan efek fase bulan.
\end{document}
